% =====================================================================
%  CTT105 - Mang may tinh | BTVN#6: Tim hieu mang WLAN | Nhom 09 - 25C04
%  Biên dịch: XeLaTeX (Overleaf: Menu > Compiler > XeLaTeX)
%  Anh: figures/ ; chu ky: figures/chuky/
% =====================================================================
\documentclass[12pt,a4paper]{article}

% ---------------- Che do ban thao ------------------------------------
% \banthaotrue : hien cac duong ke cho trong can dien.
% \banthaofalse: ban nop chinh thuc.
\newif\ifbanthao
\banthaotrue

% ---------------- Ngon ngu, font -------------------------------------
\usepackage{fontspec}
\usepackage[vietnamese]{babel}
\setmainfont{TeX Gyre Termes}
\setsansfont{TeX Gyre Heros}
\setmonofont{DejaVu Sans Mono}
\usepackage{microtype}

% ---------------- Trang, khoang cach ---------------------------------
\usepackage[a4paper,top=2.5cm,bottom=2.5cm,left=3cm,right=2cm,
headheight=15pt]{geometry}
\usepackage{setspace}
\setstretch{1.25}
\setlength{\parindent}{1.2em}
\setlength{\parskip}{0.35em}
\setlength{\emergencystretch}{3em}

% ---------------- Goi chuc nang --------------------------------------
\usepackage{xcolor}
\usepackage{graphicx}
\usepackage{amsmath}
\usepackage{booktabs,array,tabularx,longtable,multirow}
\usepackage{enumitem}
\usepackage{caption}
\usepackage{float}
\usepackage{listings}
\usepackage{tikz}
\usetikzlibrary{arrows.meta,positioning,calc,shapes.geometric,fit,backgrounds}
\usepackage{titlesec}
\usepackage{fancyhdr}
\usepackage{qrcode}
\usepackage{needspace}
\usepackage{etoolbox}
\usepackage[hidelinks]{hyperref}
\usepackage{bookmark}
\graphicspath{{./}{figures/}}

% Mau cho so do TikZ: chi dung thang xam
\definecolor{chude}{gray}{0.10}
\definecolor{nhat}{gray}{0.93}
\definecolor{vien}{gray}{0.55}
\definecolor{canhbao}{gray}{0.25}

% ---------------- Tieu de muc ----------------------------------------
\titleformat{\section}{\normalfont\large\bfseries}{\thesection.}{0.6em}{}
\titleformat{\subsection}{\normalfont\normalsize\bfseries}{\thesubsection.}{0.6em}{}
\titleformat{\subsubsection}{\normalfont\normalsize\bfseries\itshape}{\thesubsubsection.}{0.6em}{}
\titlespacing*{\section}{0pt}{1.6em}{0.8em}
\titlespacing*{\subsection}{0pt}{1.2em}{0.5em}
\titlespacing*{\subsubsection}{0pt}{0.9em}{0.3em}
\setcounter{secnumdepth}{3}
\setcounter{tocdepth}{2}
\pretocmd{\section}{\needspace{8\baselineskip}}{}{}
\pretocmd{\subsection}{\needspace{5\baselineskip}}{}{}

% ---------------- Dau trang, chan trang ------------------------------
\pagestyle{fancy}
\fancyhf{}
\fancyhead[L]{\small\itshape Đồ án WLAN -- Nhóm 09}
\fancyhead[R]{\small\itshape CTT105 -- Mạng máy tính}
\fancyfoot[C]{\small\thepage}
\renewcommand{\headrulewidth}{0.4pt}
\fancypagestyle{plain}{\fancyhf{}\fancyfoot[C]{\small\thepage}%
\renewcommand{\headrulewidth}{0pt}}

% ---------------- Chu thich, danh so ---------------------------------
\captionsetup{font=small,labelfont=bf,labelsep=colon,justification=centering,skip=6pt}
\captionsetup[table]{position=top}
\numberwithin{figure}{section}
\numberwithin{table}{section}
\addto\captionsvietnamese{%
  \renewcommand{\listfigurename}{Danh mục hình ảnh}%
  \renewcommand{\listtablename}{Danh mục bảng}%
  \renewcommand{\refname}{Tài liệu tham khảo}%
\renewcommand{\contentsname}{Mục lục}}

% ---------------- Danh sach, bang ------------------------------------
\setlist{nosep,topsep=0.3em,leftmargin=1.6em}
\setlist[enumerate]{label=\arabic*.}
\renewcommand{\arraystretch}{1.2}
\newcolumntype{L}{>{\raggedright\arraybackslash}X}
\newcolumntype{C}{>{\centering\arraybackslash}X}
\newcolumntype{P}[1]{>{\raggedright\arraybackslash}p{#1}}

% ---------------- Ma lenh --------------------------------------------
\lstset{basicstyle=\ttfamily\small,frame=lines,framesep=4pt,xleftmargin=1em,
  breaklines=true,columns=fullflexible,keepspaces=true,
commentstyle=\itshape,aboveskip=0.6em,belowskip=0.6em}

% ---------------- Macro ----------------------------------------------
\newcommand{\fillin}[1][3cm]{\ifbanthao\rule[-2pt]{#1}{0.4pt}\else\rule[-2pt]{#1}{0.4pt}\fi}
\newcommand{\kiemtra}[1]{#1}
\newcommand{\anhhoackhung}[2]{%
  \IfFileExists{figures/#2}{\includegraphics[width=#1,height=0.5\textheight,keepaspectratio]{#2}}%
  {\fbox{\parbox[c][4cm][c]{\dimexpr#1-2\fboxsep-2\fboxrule\relax}{\centering\small
[Chưa có ảnh: \texttt{figures/\detokenize{#2}}]}}}}
\newcommand{\hinh}[4][0.92\linewidth]{%
  \begin{figure}[H]\centering\anhhoackhung{#1}{#2}\caption{#3}\label{#4}
\end{figure}}
\newcommand{\chuky}[1]{\IfFileExists{figures/chuky/#1}%
  {\includegraphics[height=1.30cm,width=3.20cm,keepaspectratio]{chuky/#1}}%
{\rule{0pt}{1.15cm}}}
\newcommand{\mono}[1]{\texttt{#1}}

% =====================================================================
\begin{document}

% ============================ BIA ====================================
\begin{titlepage}
  \centering
  {\large ĐẠI HỌC QUỐC GIA THÀNH PHỐ HỒ CHÍ MINH\par}
  {\large\bfseries TRƯỜNG ĐẠI HỌC KHOA HỌC TỰ NHIÊN\par}
  {\large\bfseries KHOA CÔNG NGHỆ THÔNG TIN\par}
  \vspace{0.15cm}
  \rule{4cm}{0.5pt}\par
  \vspace{1.2cm}
  \IfFileExists{figures/logo_hcmus.png}{\includegraphics[height=3.2cm]{logo_hcmus.png}\par\vspace{1.2cm}}%
  {\IfFileExists{figures/logo_fit.png}{\includegraphics[height=3.2cm]{logo_fit.png}\par\vspace{1.2cm}}{\vspace{2.5cm}}}
  {\large BÁO CÁO ĐỒ ÁN LÝ THUYẾT\par}
  \vspace{0.2cm}
  {\large Học phần: CTT105 -- Mạng máy tính\par}
  \vspace{1.0cm}
  {\LARGE\bfseries TÌM HIỂU MẠNG WLAN\par}
  \vspace{0.3cm}
  {\large BTVN\#6 -- Mã đồ án: WLAN\par}
  \vspace{1.4cm}
  {\large\bfseries Nhóm 09 -- Lớp 25C04\par}
  \vspace{0.5cm}
  \begin{tabular}{@{}l l l@{}}
    \toprule
    \textbf{STT} & \textbf{MSSV} & \textbf{Họ và tên}\\
    \midrule
    1 & 25127076 & Nguyễn Quốc Khánh (Trưởng nhóm)\\
    2 & 25127555 & Đỗ Tấn Hoàng Vũ\\
    3 & 25127542 & Đàm Anh Tuấn\\
    4 & 21127610 & Hà Huy Hoàng\\
    5 & 25127425 & Trương Hồng Minh\\
    \bottomrule
  \end{tabular}
  \vfill
  {\large Giảng viên hướng dẫn: Huỳnh Thị Bảo Trân\par}
  \vspace{0.9cm}
  {\large Thành phố Hồ Chí Minh, tháng 11 năm 2026\par}
\end{titlepage}

\pagenumbering{roman}
\setcounter{page}{2}

% ========================= LOI CAM ON ================================
\section*{Lời cảm ơn}
\addcontentsline{toc}{section}{Lời cảm ơn}
Nhóm 09 xin gửi lời cảm ơn chân thành đến cô Huỳnh Thị Bảo Trân, giảng viên
học phần Mạng máy tính, đã hướng dẫn và định hướng nội dung để nhóm hoàn thành
đồ án \emph{Tìm hiểu mạng WLAN}. Những yêu cầu cụ thể của đồ án, từ phần lý
thuyết đến phần cấu hình thực tế có minh chứng, đã giúp nhóm hiểu rõ hơn cách
một mạng không dây gia đình được xây dựng, vận hành và bảo vệ.

Quá trình thực hiện cũng giúp các thành viên rèn luyện kỹ năng đọc tài liệu
tiếng Anh chuyên ngành, cấu hình thiết bị mạng, kiểm thử, ghi nhận minh chứng
và phối hợp làm việc nhóm. Dù đã cố gắng, báo cáo khó tránh khỏi thiếu sót;
nhóm rất mong nhận được góp ý của cô để hoàn thiện hơn.

\vspace{1em}
\hfill
\begin{tabular}{c}\textit{TP. Hồ Chí Minh, ngày 06 tháng 11 năm 2026}\\
  \textbf{Nhóm 09 -- Lớp 25C04}
\end{tabular}

% ===================== THONG TIN NHOM VA VIDEO =======================
\clearpage
\section*{Thông tin nhóm và tỷ lệ đóng góp}
\addcontentsline{toc}{section}{Thông tin nhóm và tỷ lệ đóng góp}

\begin{table}[H]
  \centering\small
  \caption*{\textbf{Bảng A.} Thông tin thành viên nhóm 09}
  \begin{tabularx}{\linewidth}{@{}c l l L@{}}
    \toprule
    \textbf{STT} & \textbf{MSSV} & \textbf{Họ và tên} & \textbf{Email}\\
    \midrule
    1 & 25127076 & Nguyễn Quốc Khánh (Trưởng nhóm) & \mono{nqkhanh2510@clc.fitus.edu.vn}\\
    2 & 25127555 & Đỗ Tấn Hoàng Vũ & \mono{dthvu2518@clc.fitus.edu.vn}\\
    3 & 25127542 & Đàm Anh Tuấn & \mono{datuan2535@clc.fitus.edu.vn}\\
    4 & 21127610 & Hà Huy Hoàng & \mono{hhhoang21@clc.fitus.edu.vn}\\
    5 & 25127425 & Trương Hồng Minh & \mono{thminh2518@clc.fitus.edu.vn}\\
    \bottomrule
  \end{tabularx}
\end{table}

\begin{table}[H]
  \centering\small
  \caption*{\textbf{Bảng B.} Phân công và tỷ lệ đóng góp (có chữ ký xác nhận)}
  \begin{tabularx}{\linewidth}{@{}c P{3.6cm} L c >{\centering\arraybackslash}p{3.4cm}@{}}
    \toprule
    \textbf{STT} & \textbf{Họ và tên} & \textbf{Tỷ lệ} & \textbf{Chữ ký}\\
    \midrule
    1 & Nguyễn Quốc Khánh & 100\,\% & \chuky{khanh.png}\\
    2 & Đỗ Tấn Hoàng Vũ   & 100\,\% & \chuky{vu.png}\\
    3 & Đàm Anh Tuấn      & 100\,\% & \chuky{tuan.png}\\
    4 & Hà Huy Hoàng      & 100\,\% & \chuky{hoang.png}\\
    5 & Trương Hồng Minh  & 80\,\% & \chuky{minh.png}\\
    \bottomrule
  \end{tabularx}
\end{table}

\subsection*{Video thực hành}
\noindent
\begin{minipage}[c]{0.76\linewidth}
  \begin{tabular}{@{}l l@{}}
    Tệp video & \mono{N09\_WLAN-video.mp4}\\
    Thời lượng & 40 phút 13 giây\\
    Dựng video & Nguyễn Quốc Khánh\\
  \end{tabular}

  \smallskip
  {\raggedright\noindent Liên kết:\\
  {\small\url{https://youtu.be/FPLgrdk_lfI}}\par}
  \noindent Liên kết cũng được lưu trong tệp \mono{N09\_WLAN-video.txt} nộp kèm. Bản đồ
  phân cảnh ở Phụ lục~\ref{app:video}.
\end{minipage}\hfill
\begin{minipage}[c]{0.2\linewidth}\centering
  \qrcode[height=2.4cm]{https://youtu.be/FPLgrdk_lfI}
\end{minipage}

% ========================= MUC LUC ===================================
\clearpage
\tableofcontents
\clearpage
\listoffigures
\listoftables
\clearpage
\pagenumbering{arabic}

% =====================================================================
\section{Giới thiệu đồ án}
\label{sec:gioithieu}

\subsection{Mục tiêu}
Đồ án tìm hiểu mạng cục bộ không dây (WLAN) và triển khai thực tế một mạng
WLAN gia đình theo hai mô hình \emph{Infrastructure} và \emph{Ad-hoc}. Cụ thể,
nhóm hướng tới bốn mục tiêu:
\begin{enumerate}
  \item Hệ thống hóa khái niệm, thành phần và chuẩn kỹ thuật của WLAN.
  \item Phân tích hai mô hình triển khai và giải pháp kết nối cho Laptop,
    Desktop và Smartphone.
  \item Cấu hình một bộ định tuyến thật thành mạng WLAN mang tên nhóm, kết nối các
    loại thiết bị và kiểm tra bằng minh chứng cụ thể.
  \item Dựng và kiểm chứng một mạng Ad-hoc tối thiểu hai nút.
\end{enumerate}

\subsection{Phạm vi}
Đồ án giới hạn trong mạng gia đình và văn phòng nhỏ: một điểm truy cập, một
mạng con IPv4 và các thiết bị đầu cuối phổ thông. Các giải pháp doanh nghiệp
(bộ điều khiển WLC, xác thực 802.1X, mạng nhiều AP) chỉ được giới thiệu ở mức
khái niệm để làm rõ vị trí của mô hình gia đình.

\subsection{Thiết bị và môi trường thực hiện}
\begin{table}[H]
  \centering\small
  \caption{Thiết bị và phần mềm sử dụng trong đồ án}
  \label{tab:thietbi}
  \begin{tabularx}{\linewidth}{@{}P{3.4cm} L@{}}
    \toprule
    \textbf{Hạng mục} & \textbf{Chi tiết}\\
    \midrule
    Router / AP & TP-Link Archer C20 v5 (AC750), firmware OpenWrt
    25.12.6, giao diện quản trị LuCI; tên thiết bị \mono{N09-Router}\\
    Desktop & Máy bàn chạy Linux Mint, có card Wi-Fi (hostname \mono{luckmintpc})\\
    Laptop & Laptop Windows (thao tác: Đỗ Tấn Hoàng Vũ)\\
    Smartphone & Điện thoại Android (thao tác: Đỗ Tấn Hoàng Vũ); Redmi Note 14 (Đàm Anh Tuấn) làm máy thử\\
    Máy thử phía WAN & Laptop Arch Linux (của Nguyễn Quốc Khánh, thao tác: Đàm Anh Tuấn), dùng để truy cập từ mạng gia đình ở các cảnh cần Internet\\
    Mô hình Ad-hoc & Máy ảo Ubuntu 22.04 LTS trên VMware, module
    \mono{mac80211\_hwsim} (thao tác: Hà Huy Hoàng)\\
    Công cụ & Trình duyệt web, \mono{nmcli}, \mono{iw}, \mono{ip}, \mono{ping},
    \mono{netsh}, \mono{ipconfig}; OBS Studio để quay màn hình\\
    \bottomrule
  \end{tabularx}
\end{table}

\subsection{Cấu trúc báo cáo}
Mục~\ref{sec:khainiem}--\ref{sec:baomat} trình bày lý thuyết: khái niệm,
thành phần, hai mô hình và bảo mật. Mục~\ref{sec:giaiphap} là giải pháp kết
nối cho từng loại thiết bị. Mục~\ref{sec:infra} và \ref{sec:adhoc} là phần
thực nghiệm; Mục~\ref{sec:ketqua} tổng hợp kết quả, đối chiếu với yêu cầu đề
và thảo luận.

% =====================================================================
\section{Khái niệm và tiêu chuẩn WLAN}
\label{sec:khainiem}

\subsection{Khái niệm WLAN}
\textbf{WLAN} (Wireless Local Area Network) là mạng cục bộ trong đó các thiết
bị trao đổi dữ liệu bằng sóng vô tuyến thay vì cáp. Hầu hết WLAN hiện nay xây
dựng trên họ chuẩn \textbf{IEEE~802.11}, với tên thương mại \emph{Wi-Fi} do
Wi-Fi Alliance chứng nhận \cite{ieee80211,gast2005}. Về phạm vi, WLAN nằm giữa
mạng cá nhân không dây (WPAN, vài mét, ví dụ Bluetooth) và mạng diện rộng không
dây (WWAN, hàng kilômét, ví dụ mạng di động 4G/5G).

Khác với Ethernet, môi trường truyền của WLAN là không gian dùng chung: mọi
thiết bị trong vùng phủ cùng nghe một kênh. Từ đặc điểm này sinh ra cả ưu
điểm lẫn hạn chế của WLAN (Bảng~\ref{tab:uunhuoc}).

\begin{table}[H]
  \centering\small
  \caption{Ưu điểm và hạn chế của WLAN so với LAN có dây}
  \label{tab:uunhuoc}
  \begin{tabularx}{\linewidth}{@{}L L@{}}
    \toprule
    \textbf{Ưu điểm} & \textbf{Hạn chế}\\
    \midrule
    Thiết bị di chuyển tự do trong vùng phủ sóng &
    Băng thông chia sẻ giữa mọi thiết bị trên cùng kênh\\
    Triển khai nhanh, không cần đi dây tới từng vị trí &
    Chất lượng phụ thuộc khoảng cách, vật cản và nhiễu\\
    Dễ mở rộng bằng cách thêm điểm truy cập &
    Độ trễ cao và dao động hơn mạng có dây\\
    Phù hợp với thiết bị di động và IoT &
    Sóng lan ra ngoài nên bắt buộc phải mã hóa và xác thực\\
    \bottomrule
  \end{tabularx}
\end{table}

\subsection{Các thế hệ chuẩn IEEE 802.11}
Chuẩn 802.11 gốc (1997) chỉ đạt 2~Mb/s. Các bản sửa đổi sau đó tăng tốc độ
nhờ điều chế bậc cao hơn, kênh rộng hơn và nhiều luồng không gian (MIMO).
Bảng~\ref{tab:thehe} tóm tắt các thế hệ chính; tốc độ trong bảng là tốc độ
lớp vật lý lý thuyết tối đa, thông lượng thực tế thấp hơn nhiều
\cite{perahia2013}.

\begin{table}[H]
  \centering\small
  \caption{Các thế hệ chuẩn IEEE 802.11}
  \label{tab:thehe}
  \begin{tabularx}{\linewidth}{@{}l c c c c L@{}}
    \toprule
    \textbf{Chuẩn} & \textbf{Tên gọi} & \textbf{Năm} & \textbf{Băng tần (GHz)} &
    \textbf{Tốc độ tối đa} & \textbf{Công nghệ chính}\\
    \midrule
    802.11   & --         & 1997 & 2,4         & 2 Mb/s    & DSSS, FHSS, CSMA/CA\\
    802.11b  & --         & 1999 & 2,4         & 11 Mb/s   & DSSS/CCK\\
    802.11a  & --         & 1999 & 5           & 54 Mb/s   & OFDM\\
    802.11g  & --         & 2003 & 2,4         & 54 Mb/s   & OFDM, tương thích 802.11b\\
    802.11n  & Wi-Fi 4    & 2009 & 2,4 / 5     & 600 Mb/s  & MIMO 4 luồng, kênh 40 MHz\\
    802.11ac & Wi-Fi 5    & 2013 & 5           & 6,9 Gb/s  & 256-QAM, kênh 160 MHz, MU-MIMO đường xuống\\
    802.11ax & Wi-Fi 6/6E & 2021 & 2,4 / 5 / 6 & 9,6 Gb/s  & OFDMA, 1024-QAM, TWT\\
    802.11be & Wi-Fi 7    & 2024 & 2,4 / 5 / 6 & $\approx$46 Gb/s & 4096-QAM, kênh 320 MHz, MLO\\
    \bottomrule
  \end{tabularx}
\end{table}

Ngoài các bản sửa đổi lớp vật lý, nhiều bản sửa đổi bổ sung chức năng cho lớp
MAC: 802.11e (chất lượng dịch vụ), 802.11i và 802.11w (bảo mật, bảo vệ khung
quản lý), 802.11k/v/r (hỗ trợ chuyển vùng) và 802.11s (mạng mesh).

Router dùng trong đồ án là loại AC750: băng 2,4~GHz theo 802.11n (tối đa
300~Mb/s) và băng 5~GHz theo 802.11ac (tối đa 433~Mb/s).

\subsection{Băng tần và kênh}
\begin{itemize}
  \item \textbf{2,4~GHz} (2400--2483,5 MHz): tầm phủ xa, xuyên tường tốt nhưng
    trong quy hoạch thường dùng, các kênh 20~MHz 1, 6 và 11 được xem là không
    chồng lấn; số kênh và mức chồng lấn thực tế phụ thuộc khu vực pháp lý,
    băng thông và quy hoạch kênh. Việt Nam cho phép kênh 1--13.
  \item \textbf{5~GHz}: nhiều kênh hơn, ít nhiễu, ghép được kênh 40/80/160~MHz.
    Một số kênh phải dùng cơ chế DFS để nhường kênh khi phát hiện radar.
  \item \textbf{6~GHz}: dùng cho Wi-Fi 6E/7. Tại Việt Nam, Thông tư
    01/2025/TT-BKHCN cho phép dải 5925--6425~MHz \cite{tt012025}; theo yêu cầu
    chứng nhận của Wi-Fi Alliance, mạng ở băng này phải dùng WPA3 hoặc OWE.
\end{itemize}
Thiết bị phát Wi-Fi lưu hành tại Việt Nam phải tuân thủ quy chuẩn kỹ thuật quốc
gia (QCVN) về tần số và công suất. Vì vậy khi cấu hình router, nhóm đặt mã quốc
gia \mono{VN} để router chỉ dùng kênh và công suất được phép.

\subsection{Vị trí của WLAN trong mô hình OSI}
IEEE 802.11 chỉ định nghĩa hai tầng thấp nhất \cite{kurose2012}:
\begin{itemize}
  \item \textbf{Tầng vật lý (PHY)}: điều chế (DSSS, OFDM, QAM), tần số, kênh
    và công suất phát.
  \item \textbf{Tầng liên kết dữ liệu}: gồm phân tầng \textbf{MAC 802.11}
    (truy cập kênh theo CSMA/CA, xác nhận ACK, định địa chỉ khung, mã hóa) và
    phân tầng \textbf{LLC 802.2} che đi khác biệt giữa mạng có dây và không
    dây đối với tầng mạng.
\end{itemize}
Từ tầng mạng trở lên (IP, TCP/UDP, ứng dụng), một thiết bị Wi-Fi hoạt động y
hệt thiết bị cắm dây. Do đó WLAN vẫn cần các dịch vụ IP như DHCP và DNS
(Mục~\ref{subsec:dichvu}).

Điểm khác biệt quan trọng của tầng MAC là cơ chế \textbf{CSMA/CA}
(\emph{tránh} xung đột) thay cho CSMA/CD của Ethernet: thiết bị vô tuyến không
thể vừa phát vừa nghe để phát hiện xung đột. Trước khi phát, trạm lắng nghe
kênh và chờ một khoảng ngẫu nhiên (backoff); mỗi khung unicast phải được bên
nhận trả về ACK.

% =====================================================================
\section{Thành phần của mạng WLAN}
\label{sec:thanhphan}

\subsection{Thành phần phần cứng}
\begin{table}[H]
  \centering\small
  \caption{Các thành phần phần cứng của WLAN}
  \label{tab:thanhphan}
  \begin{tabularx}{\linewidth}{@{}P{3.8cm} L@{}}
    \toprule
    \textbf{Thành phần} & \textbf{Vai trò}\\
    \midrule
    Trạm (Station -- STA) & Thiết bị đầu cuối có giao tiếp 802.11: laptop, điện
    thoại, máy bàn có card Wi-Fi, camera, thiết bị IoT\\
    Card mạng không dây (WNIC) & Bộ thu phát vô tuyến của trạm, gắn liền bo mạch
    hoặc rời (PCIe, USB); mang địa chỉ MAC của trạm\\
    Điểm truy cập (Access Point -- AP) & Phát sóng, quản lý các trạm liên kết và
    làm cầu nối giữa mạng không dây với mạng có dây\\
    Bộ định tuyến không dây & Thiết bị gia đình tích hợp AP, switch, bộ định tuyến, máy
    chủ DHCP, NAT và tường lửa trong một hộp\\
    Hệ thống phân phối (DS) & Hạ tầng nối các AP với nhau và với mạng ngoài,
    thường là Ethernet và switch\\
    Bộ điều khiển WLAN (WLC) & Quản lý tập trung nhiều AP trong doanh nghiệp
    (cấu hình, kênh, công suất, chuyển vùng)\\
    Ăng-ten & Bức xạ và thu sóng; quyết định hình dạng vùng phủ\\
    \bottomrule
  \end{tabularx}
\end{table}

Trong mạng gia đình, toàn bộ AP, switch, DS và router thường nằm trong một
thiết bị duy nhất. Router TP-Link Archer C20 của nhóm là ví dụ: một hộp có hai
bộ phát vô tuyến (2,4 và 5~GHz), bốn cổng LAN, một cổng WAN và phần mềm định
tuyến OpenWrt.

\subsection{Thành phần định danh}
\begin{itemize}
  \item \textbf{SSID} (Service Set Identifier): tên mạng dài tối đa 32 byte,
    được AP quảng bá trong khung beacon để người dùng chọn mạng.
  \item \textbf{BSSID}: định danh 48 bit của một BSS. Trong mô hình
    Infrastructure, BSSID chính là địa chỉ MAC của giao diện vô tuyến trên AP;
    một router hai băng tần có hai BSSID khác nhau dù có thể cùng SSID.
  \item \textbf{Địa chỉ MAC}: địa chỉ 48 bit của card mạng. Bit thấp thứ hai (U/L) của
    byte đầu tiên cho biết địa chỉ do nhà sản xuất cấp (bit = 0) hay được gán
    cục bộ (bit = 1). Các hệ điều hành hiện đại dùng \emph{địa chỉ MAC ngẫu
    nhiên} (gán cục bộ) cho từng mạng Wi-Fi để bảo vệ quyền riêng tư; điều này
    ảnh hưởng trực tiếp tới các chức năng lọc theo MAC của router
    (Mục~\ref{subsec:macngaunhien}).
\end{itemize}

\subsection{Dịch vụ IP đi kèm}
\label{subsec:dichvu}
\begin{itemize}
  \item \textbf{DHCP} cấp tự động địa chỉ IP, mặt nạ mạng, cổng mặc định và DNS
    cho trạm qua bốn bước Discover -- Offer -- Request -- Acknowledge (DORA).
  \item \textbf{DNS} phân giải tên miền thành địa chỉ IP. Router gia đình
    thường làm máy chủ DNS chuyển tiếp cho mạng nội bộ.
  \item \textbf{Cổng mặc định (default gateway)} là địa chỉ LAN của router,
    nơi trạm gửi mọi gói tin đi ra ngoài mạng con.
  \item \textbf{NAT} cho phép mọi thiết bị dùng địa chỉ riêng (ví dụ
    192.168.1.0/24) chia sẻ một địa chỉ phía WAN để ra Internet.
\end{itemize}

\subsection{Quá trình một trạm gia nhập mạng}
Hình~\ref{fig:ketnoi} mô tả các bước một trạm thực hiện để vào được mạng
WPA2-Personal và có địa chỉ IP. Đây cũng là trình tự cần kiểm tra khi kết nối
thất bại (Mục~\ref{subsec:loi}).

\begin{figure}[H]
  \centering
  \begin{tikzpicture}[font=\small,
      msg/.style={-{Stealth[length=2mm]},thick,chude},
      lbl/.style={fill=white,inner sep=1.5pt,font=\footnotesize},
    phase/.style={font=\footnotesize\bfseries,text=canhbao,anchor=west}]
    \node[draw=chude,fill=nhat,minimum width=2.6cm,minimum height=0.7cm] (sta) at (0,0) {Trạm (STA)};
    \node[draw=chude,fill=nhat,minimum width=2.6cm,minimum height=0.7cm] (ap) at (7,0) {AP / Router};
    \draw[vien,dashed] (sta.south) -- ++(0,-10.2);
    \draw[vien,dashed] (ap.south) -- ++(0,-10.2);
    \draw[msg,dashed] (7,-0.9) -- node[lbl]{Beacon (SSID, khả năng, bảo mật)} (0,-0.9);
    \draw[msg] (0,-1.6) -- node[lbl]{Probe Request} (7,-1.6);
    \draw[msg] (7,-2.2) -- node[lbl]{Probe Response} (0,-2.2);
    \draw[msg] (0,-3.0) -- node[lbl]{Authentication (Open System)} (7,-3.0);
    \draw[msg] (7,-3.6) -- node[lbl]{Authentication Response} (0,-3.6);
    \draw[msg] (0,-4.4) -- node[lbl]{Association Request} (7,-4.4);
    \draw[msg] (7,-5.0) -- node[lbl]{Association Response (AID)} (0,-5.0);
    \draw[msg] (7,-5.8) -- node[lbl]{EAPOL-Key 1/4 (ANonce)} (0,-5.8);
    \draw[msg] (0,-6.35) -- node[lbl]{EAPOL-Key 2/4 (SNonce, MIC)} (7,-6.35);
    \draw[msg] (7,-6.9) -- node[lbl]{EAPOL-Key 3/4 (GTK, MIC)} (0,-6.9);
    \draw[msg] (0,-7.45) -- node[lbl]{EAPOL-Key 4/4 (xác nhận cài khóa)} (7,-7.45);
    \draw[msg] (0,-8.3) -- node[lbl]{DHCP Discover} (7,-8.3);
    \draw[msg] (7,-8.85) -- node[lbl]{DHCP Offer} (0,-8.85);
    \draw[msg] (0,-9.4) -- node[lbl]{DHCP Request} (7,-9.4);
    \draw[msg] (7,-9.95) -- node[lbl]{DHCP ACK (IP, gateway, DNS)} (0,-9.95);
    \node[phase] at (7.3,-1.9) {Dò tìm};
    \node[phase] at (7.3,-3.3) {Xác thực 802.11};
    \node[phase] at (7.3,-4.7) {Liên kết};
    \node[phase] at (7.3,-6.62) {Bắt tay 4 bước (WPA2)};
    \node[phase] at (7.3,-9.12) {Cấp địa chỉ IP (DORA)};
  \end{tikzpicture}
  \caption{Trình tự một trạm gia nhập mạng WPA2-Personal và nhận địa chỉ IP}
  \label{fig:ketnoi}
\end{figure}

% =====================================================================
\section{Mô hình triển khai: Infrastructure và Ad-hoc}
\label{sec:mohinh}

Chuẩn 802.11 tổ chức mạng thành các \emph{tập dịch vụ} (service set). Có hai
cách tổ chức cơ bản: có điểm truy cập trung tâm (Infrastructure) và không có
điểm truy cập (Ad-hoc).

\subsection{Mô hình Infrastructure}
\subsubsection{Khái niệm và thành phần}
Trong mô hình Infrastructure, mọi trạm liên kết với một AP và mọi khung dữ
liệu đều đi qua AP, kể cả khi hai trạm nằm cạnh nhau. Mô hình gồm ba nhóm
thành phần: AP hoặc router không dây làm trung tâm; các trạm đầu cuối; và hạ
tầng có dây (DS) nối AP với mạng ngoài. Hình~\ref{fig:thanhphan} minh họa các
thành phần này trong một mạng gia đình điển hình.

\hinh[\linewidth]{N09_sodo3.pdf}{Các thành phần của mạng WLAN gia đình theo mô hình Infrastructure: một router tích hợp tạo BSS 1, AP thứ hai nối qua hệ thống phân phối tạo BSS 2, hai BSS cùng SSID hợp thành một ESS}{fig:thanhphan}

\subsubsection{BSS, ESS và chuyển vùng}
\begin{itemize}
  \item \textbf{BSS} (Basic Service Set) là đơn vị cơ bản: một AP cùng các trạm
    đang liên kết với nó, định danh bằng BSSID.
  \item \textbf{ESS} (Extended Service Set) gồm nhiều BSS nối qua DS, dùng
    chung SSID và cấu hình bảo mật để phủ sóng diện rộng. Mỗi AP vẫn giữ BSSID
    riêng.
  \item \textbf{Chuyển vùng (roaming)}: trong ESS, trạm tự quyết định chuyển
    sang AP có tín hiệu tốt hơn mà vẫn giữ địa chỉ IP nếu cùng mạng con. Các
    bản sửa đổi 802.11k/v/r giúp quá trình này nhanh và chính xác hơn.
\end{itemize}

\subsection{Mô hình Ad-hoc (IBSS)}
\subsubsection{Khái niệm và thành phần}
Mô hình Ad-hoc, trong chuẩn gọi là \textbf{IBSS} (Independent BSS), gồm các
trạm liên lạc trực tiếp với nhau, không có AP. Thành phần chỉ là các trạm có
card Wi-Fi hỗ trợ chế độ IBSS, cùng cấu hình chung \textbf{SSID} và
\textbf{kênh}. Các đặc điểm kỹ thuật chính:
\begin{itemize}
  \item \textbf{BSSID ngẫu nhiên}: trạm khởi tạo mạng sinh một BSSID ngẫu
    nhiên dạng gán cục bộ; các trạm gia nhập sau dùng chung BSSID này.
  \item \textbf{Beacon phân tán}: không có AP nên các trạm cùng chịu trách
    nhiệm phát beacon. Sau mỗi chu kỳ beacon, mỗi trạm chờ một khoảng ngẫu
    nhiên; trạm nào phát trước thì các trạm khác hủy lượt phát của mình.
  \item \textbf{Không có dịch vụ tập trung}: không có DHCP nên địa chỉ IP phải
    gán tĩnh hoặc tự cấu hình (link-local). Mạng mặc định cô lập, không có
    Internet trừ khi một nút được cấu hình thêm làm cổng ra.
  \item \textbf{Một chặng}: IBSS cơ bản chỉ cho các trạm trong tầm sóng của
    nhau liên lạc trực tiếp; trạm \emph{không} tự chuyển tiếp gói cho trạm
    khác. Muốn truyền qua nhiều chặng phải bổ sung giao thức định tuyến
    (MANET với AODV, OLSR, B.A.T.M.A.N.) hoặc dùng mesh 802.11s.
\end{itemize}

\subsubsection{Tình trạng hỗ trợ hiện nay}
IBSS hỗ trợ bảo mật kém hơn mô hình Infrastructure: giao diện của nhiều hệ
điều hành cũ chỉ cho WEP hoặc mạng mở, dù trên Linux có thể chạy IBSS với
WPA2-PSK thông qua \mono{wpa\_supplicant}. Khả năng tạo IBSS trên hệ điều hành phổ thông hiện không đồng nhất, phụ thuộc
phiên bản, driver và phần cứng; Linux thường cung cấp công cụ cấu hình trực
tiếp đầy đủ hơn. Android và iOS thường không cung cấp giao diện người dùng để
tạo IBSS. Nhu cầu kết nối trực tiếp giữa hai thiết bị nay thường được đáp ứng
bằng \textbf{Wi-Fi Direct}. Cần lưu ý Wi-Fi
Direct \emph{không phải} IBSS: một thiết bị được bầu làm \emph{Group Owner} và
hoạt động như một AP mềm, còn các thiết bị khác liên kết vào như trạm thường.

Ad-hoc vẫn có giá trị trong các mạng chuyên dụng không có hạ tầng: MANET (cứu
hộ, quân sự), VANET (liên lạc giữa xe), FANET (bầy máy bay không người lái) và
mạng cảm biến. Các mạng này dựa trên IBSS hoặc mesh, cộng thêm giao thức định
tuyến nhiều chặng.

\subsection{So sánh và lựa chọn mô hình}
\begin{table}[H]
  \centering\small
  \caption{So sánh mô hình Infrastructure và Ad-hoc}
  \label{tab:sosanh}
  \begin{tabularx}{\linewidth}{@{}P{3.1cm} L L@{}}
    \toprule
    \textbf{Tiêu chí} & \textbf{Infrastructure (BSS/ESS)} & \textbf{Ad-hoc (IBSS)}\\
    \midrule
    Thiết bị trung tâm & Bắt buộc có AP/router & Không có\\
    Đường đi dữ liệu & Mọi khung đi qua AP & Trực tiếp giữa hai trạm\\
    Định danh BSS & BSSID = MAC vô tuyến của AP & BSSID ngẫu nhiên do trạm đầu tiên sinh\\
    Beacon & Do AP phát & Các trạm luân phiên phát\\
    Cấp địa chỉ IP & DHCP trên router & Gán tĩnh hoặc link-local\\
    Kết nối Internet & Có, qua router và NAT & Mặc định không có\\
    Bảo mật & WPA2/WPA3, 802.1X, quản lý tập trung & Hạn chế, phụ thuộc hệ điều hành và driver\\
    Mở rộng & Thêm AP thành ESS, có chuyển vùng & Kém, hiệu năng giảm khi nhiều nút\\
    Hỗ trợ trên hệ điều hành & Mọi hệ điều hành & Phụ thuộc hệ điều hành, driver và phần cứng\\
    Ứng dụng & Gia đình, văn phòng, trường học & Kết nối tạm thời, MANET, cảm biến\\
    \bottomrule
  \end{tabularx}
\end{table}

Với mạng gia đình, Infrastructure là lựa chọn phù hợp vì có Internet, cấp
địa chỉ tự động, bảo mật hiện đại và quản lý tập trung. Ad-hoc chỉ phù hợp khi
không có hạ tầng hoặc để nghiên cứu nguyên lý kết nối ngang hàng, đúng như
cách nhóm sử dụng nó trong Mục~\ref{sec:adhoc}.

% =====================================================================
\section{Bảo mật WLAN}
\label{sec:baomat}
Sóng vô tuyến lan ra ngoài phạm vi ngôi nhà, nên bất kỳ ai trong vùng phủ đều
thu được khung tin. Bảo mật WLAN vì vậy phải bảo đảm cả xác thực (chỉ người có
quyền mới vào được) và mã hóa (người nghe lén không đọc được dữ liệu).

\begin{table}[H]
  \centering\small
  \caption{Các thế hệ cơ chế bảo mật Wi-Fi}
  \label{tab:baomat}
  \begin{tabularx}{\linewidth}{@{}l P{3.6cm} L@{}}
    \toprule
    \textbf{Cơ chế} & \textbf{Mã hóa / xác thực} & \textbf{Đánh giá}\\
    \midrule
    WEP & RC4, khóa tĩnh & Đã bị phá hoàn toàn \cite{fms2001}; không được dùng\\
    WPA & TKIP & Giải pháp tạm thời, đã lỗi thời\\
    WPA2 & AES-CCMP, bắt tay 4 bước & Phổ biến nhất; chế độ PSK có thể bị dò mật
    khẩu ngoại tuyến nếu mật khẩu yếu; lỗ hổng KRACK \cite{krack2017} đã được vá\\
    WPA3 & AES-CCMP/GCMP, SAE & Chống dò mật khẩu ngoại tuyến, bắt buộc bảo vệ
    khung quản lý (PMF) \cite{wpa3spec}\\
    OWE & Mã hóa không cần mật khẩu & Dành cho mạng công cộng mở (Wi-Fi Enhanced Open)\\
    \bottomrule
  \end{tabularx}
\end{table}

Mỗi chuẩn WPA có hai chế độ xác thực:
\begin{itemize}
  \item \textbf{Personal (PSK/SAE)}: mọi người dùng chung một mật khẩu; phù hợp
    gia đình.
  \item \textbf{Enterprise (802.1X)}: mỗi người dùng xác thực riêng qua EAP
    với máy chủ RADIUS; phù hợp doanh nghiệp. Chế độ này \emph{bắt buộc} phải
    có máy chủ RADIUS. Nếu chọn nhầm chế độ Enterprise trên router gia đình,
    SSID nhiều khả năng sẽ không hoạt động được (một tình huống nhóm gặp thực tế, xem
    Mục~\ref{subsec:nhanxet}).
\end{itemize}

Các mối đe dọa thường gặp gồm nghe lén, dò mật khẩu PSK, AP giả mạo (rogue
AP, evil twin), tấn công hủy xác thực (deauthentication) để ngắt kết nối, và
khai thác lỗ hổng WPS. Với mạng gia đình, nhóm đề xuất các biện pháp tối thiểu:
\begin{enumerate}
  \item Dùng WPA3 hoặc WPA2 với AES-CCMP; tắt WEP, TKIP và WPS.
  \item Đặt mật khẩu Wi-Fi dài, ngẫu nhiên; đổi mật khẩu quản trị router ngay
    lần đăng nhập đầu tiên.
  \item Bật bảo vệ khung quản lý (PMF/802.11w) nếu thiết bị hỗ trợ.
  \item Tách mạng khách khỏi mạng chính; cập nhật firmware định kỳ.
  \item Không xem ẩn SSID hay lọc MAC là biện pháp bảo mật chính: SSID vẫn lộ
    trong khung dò tìm và địa chỉ MAC có thể bị giả mạo.
\end{enumerate}

% =====================================================================
\section{Giải pháp kết nối và cấu hình thiết bị}
\label{sec:giaiphap}

\subsection{Yêu cầu chung}
Để một thiết bị vào được mạng Infrastructure cần bốn điều kiện: (1) có card
Wi-Fi và driver hoạt động; (2) card hỗ trợ băng tần và chuẩn bảo mật mà AP
dùng (ví dụ card chỉ có 2,4~GHz sẽ không thấy mạng 5~GHz); (3) biết đúng SSID
và mật khẩu; (4) nhận được cấu hình IP, thường qua DHCP. Sau khi kết nối, mọi
thiết bị cần được kiểm tra theo cùng một trình tự: địa chỉ IP thuộc đúng mạng
con, cổng mặc định là địa chỉ router, ping tới router thành công, rồi ping tới
thiết bị khác hoặc ra Internet.

Máy bàn (Desktop) thường không có sẵn Wi-Fi. Có ba lựa chọn: lắp card Wi-Fi
PCIe, dùng USB Wi-Fi, hoặc nối cáp Ethernet vào cổng LAN của router. Trong đồ
án, máy bàn của nhóm có card Wi-Fi nên được kết nối bằng Wi-Fi để đúng yêu cầu
đề bài.

\subsection{Thao tác theo từng hệ điều hành}
\begin{table}[H]
  \centering\small
  \caption{Thao tác kết nối và kiểm tra Wi-Fi trên từng hệ điều hành}
  \label{tab:hedieuhanh}
  \begin{tabularx}{\linewidth}{@{}P{2.0cm} L L L@{}}
    \toprule
    \textbf{Hệ điều hành} & \textbf{Kiểm tra card} & \textbf{Kết nối} &
    \textbf{Xem IP, MAC, gateway}\\
    \midrule
    Windows 10/11 & \mono{netsh wlan show drivers} (chuẩn, băng tần hỗ trợ) &
    Biểu tượng Wi-Fi trên thanh tác vụ, chọn SSID, nhập mật khẩu &
    \mono{ipconfig /all}; \mono{netsh wlan show interfaces} (SSID, BSSID, kênh,
    kiểu bảo mật)\\
    Linux (Mint, Ubuntu) & \mono{ip link}, \mono{iw dev}, \mono{iw list} &
    Applet mạng hoặc \mono{nmcli device wifi connect <SSID> password <...>} &
    \mono{ip addr}, \mono{ip route}, \mono{iw dev <card> link}\\
    Android & Cài đặt $\rightarrow$ Wi-Fi &
    Chọn SSID; mục \emph{Quyền riêng tư} chọn \emph{Dùng MAC của thiết bị} &
    Chạm vào mạng đang kết nối: IP, gateway, tần số, tốc độ liên kết\\
    iOS & Cài đặt $\rightarrow$ Wi-Fi &
    Chọn SSID; trong nút (i) tắt \emph{Địa chỉ Wi-Fi riêng tư} &
    Nút (i): địa chỉ IP, mặt nạ mạng, bộ định tuyến\\
    \bottomrule
  \end{tabularx}
\end{table}

\subsection{Địa chỉ MAC ngẫu nhiên}
\label{subsec:macngaunhien}
Android 10 trở lên và iOS 14 trở lên mặc định dùng một địa chỉ
MAC ngẫu nhiên, khác nhau cho mỗi mạng Wi-Fi; Windows 10/11 có tùy chọn tương tự
(\emph{Random hardware addresses}, mặc định tắt). Khi đó địa chỉ router nhìn thấy
không phải MAC thật của máy, nên các chức năng dựa trên MAC (lọc MAC, cấp IP cố
định, chặn theo giờ) sẽ áp dụng sai thiết bị hoặc mất tác dụng khi địa chỉ đổi.
Có thể nhận ra MAC ngẫu nhiên nhờ chữ số thứ hai của địa chỉ là 2, 6, A hoặc E
(bit gán cục bộ bằng 1). Vì vậy, với mạng của nhóm, các thiết bị phải tắt tính
năng này trước khi thử các chức năng quản lý.

\subsection{Lỗi thường gặp và cách khắc phục}
\label{subsec:loi}
Bảng~\ref{tab:loi} tổng hợp các lỗi thường gặp khi kết nối WLAN, trong đó các
dòng đánh dấu $\star$ là tình huống nhóm gặp trực tiếp khi thực hiện đồ án.

\begin{longtable}{@{}P{3.3cm} P{5.4cm} P{6.0cm}@{}}
  \caption{Lỗi thường gặp khi kết nối WLAN}\label{tab:loi}\\
  \toprule
  \textbf{Hiện tượng} & \textbf{Nguyên nhân} & \textbf{Cách khắc phục}\\
  \midrule\endfirsthead
  \toprule
  \textbf{Hiện tượng} & \textbf{Nguyên nhân} & \textbf{Cách khắc phục}\\
  \midrule\endhead
  \bottomrule\endfoot
  Không thấy SSID trong danh sách & Card không hỗ trợ băng tần của AP; SSID bị
  ẩn; bộ phát đang tắt & Kiểm tra băng tần card hỗ trợ; nhập SSID thủ công;
  bật bộ phát trên router\\
  $\star$ SSID không phát dù đã bật & Chọn nhầm kiểu mã hóa \emph{WPA2-EAP
  (802.1X)} khi không có máy chủ RADIUS & Chọn lại \emph{WPA2-PSK} hoặc
  \emph{WPA3-SAE}, áp dụng lại cấu hình\\
  Báo sai mật khẩu & Gõ sai, phân biệt hoa thường; thiết bị cũ không hỗ trợ WPA3
  & Nhập lại; dùng chế độ hỗn hợp WPA2/WPA3 cho thiết bị cũ\\
  Kết nối được nhưng IP dạng 169.254.x.x & Không nhận được DHCP (DHCP tắt, hết dải
  địa chỉ) & Bật hoặc mở rộng dải DHCP; ngắt và kết nối lại\\
  Bị từ chối dù đúng mật khẩu & MAC của thiết bị nằm trong danh sách lọc & Gỡ
  MAC khỏi danh sách hoặc tắt lọc MAC\\
  $\star$ Luật theo MAC không có tác dụng & Thiết bị dùng MAC ngẫu nhiên & Tắt
  MAC ngẫu nhiên cho mạng đó; cập nhật lại MAC trong luật\\
  $\star$ Có Wi-Fi nhưng không vào được web & Cổng WAN chưa nối hoặc chưa có địa
  chỉ; DNS lỗi & Kiểm tra trạng thái WAN trên router; ping IP để tách lỗi DNS\\
  Ping tới laptop Windows thất bại & Tường lửa Windows chặn ICMP đến & Thêm luật
  cho phép ICMPv4 Echo Request hoặc kiểm tra bằng chiều ngược lại\\
  $\star$ Ping tới thiết bị báo \emph{Destination Host Unreachable} & Thiết bị
  đã ngắt hoặc card Wi-Fi đang ngủ, không trả lời ARP; bản ghi DHCP vẫn còn &
  Mở màn hình thiết bị; đối chiếu danh sách \emph{Associated Stations}\\
  $\star$ Mất kết nối và cấu hình tự quay lại sau khi Apply & LuCI tự hoàn tác
  sau 90 giây khi máy quản trị mất kết nối (đổi SSID, đổi IP) & Dùng
  \emph{Apply unchecked} hoặc quản trị qua cáp LAN\\
  Không vào được trang quản trị & Sai địa chỉ gateway; đang ở mạng khách bị
  chặn truy cập router & Xem gateway bằng \mono{ipconfig}/\mono{ip route};
  chuyển sang mạng chính\\
  Router không ra Internet khi đặt sau router khác & Hai router dùng trùng dải
  địa chỉ LAN & Đổi dải LAN của một trong hai router (ví dụ 192.168.10.0/24)\\
\end{longtable}

% =====================================================================
\section{Thực nghiệm mô hình Infrastructure}
\label{sec:infra}

\subsection{Mô hình và thông số mạng}
Mạng của nhóm là một BSS kép: router N09-Router phát hai SSID trên hai bộ phát
độc lập, các trạm liên kết trực tiếp vào router (Hình~\ref{fig:sodo-infra}).
Cổng WAN được nối vào mạng gia đình (router mesh) trong hai giai đoạn: khi
thiết lập ban đầu để kiểm tra Internet, và khi thử các chức năng cần Internet
(chặn tên miền, giới hạn giờ, chuyển tiếp cổng, traceroute). Ở phần kết nối
thiết bị, cổng WAN được rút ra nên các phép thử tập trung vào liên lạc trong
mạng nội bộ.

\hinh[\linewidth]{N09_sodo1.pdf}{Sơ đồ mạng Infrastructure thực tế của nhóm}{fig:sodo-infra}

\begin{table}[H]
  \centering\small
  \caption{Thông số cấu hình thực tế của N09-Router}
  \label{tab:thongso}
  \begin{tabularx}{\linewidth}{@{}P{3.6cm} L@{}}
    \toprule
    \textbf{Hạng mục} & \textbf{Giá trị}\\
    \midrule
    Thiết bị & TP-Link Archer C20 v5, OpenWrt; hostname \mono{N09-Router},
    múi giờ Asia/Ho\_Chi\_Minh\\
    Mạng LAN (\mono{br-lan}) & \mono{192.168.1.1/24}, MAC \mono{B0:95:75:9F:1C:A3}\\
    DHCP & Dải \mono{192.168.1.100}--\mono{.249}, thời hạn thuê 12 giờ (mặc định
    OpenWrt); IP cố định \mono{192.168.1.50} cho máy bàn\\
    WAN (\mono{eth0.2}) & DHCP client, MAC \mono{B0:95:75:9F:1C:A4}\\
    Bộ phát 2,4 GHz (radio0) & MediaTek MT7628, 802.11b/g/n, kênh 6, rộng
    20~MHz, mã quốc gia VN, 22~dBm; SSID \mono{25CLC04-N09}, BSSID
    \mono{B0:95:75:9F:1C:A3}\\
    Bộ phát 5 GHz (radio1) & MediaTek MT7610E, 802.11ac, kênh 36 (5,180~GHz),
    rộng 80~MHz, mã quốc gia VN, 19~dBm; SSID \mono{25CLC04-N09-5G}, BSSID
    \mono{B0:95:75:9F:1C:A2}\\
    Bảo mật & WPA2-PSK, mã hóa CCMP (AES)\\
    \bottomrule
  \end{tabularx}
\end{table}

\subsection{Cấu hình router}
\label{subsec:cauhinh}
Router được đưa về cấu hình mặc định của OpenWrt trước buổi quay, sau đó cấu
hình từ đầu trên máy bàn nối cáp vào cổng LAN. Nhóm dùng cáp cho máy quản trị
vì khi đổi SSID hoặc kênh, các máy dùng Wi-Fi sẽ bị ngắt.

\begin{enumerate}
  \item \textbf{Đăng nhập và đặt mật khẩu quản trị.} Truy cập
    \mono{http://192.168.1.1}, đăng nhập tài khoản \mono{root} chưa có mật
    khẩu. LuCI hiện cảnh báo chưa đặt mật khẩu; nhóm đặt mật khẩu ngay tại
    \emph{System $\rightarrow$ Administration}. Trang \emph{Network
    $\rightarrow$ Wireless} lúc này cho thấy cả hai bộ phát đều tắt, đây là
    mặc định của OpenWrt để không phát ra mạng không mật khẩu.
  \item \textbf{Đặt tên thiết bị và múi giờ.} Tại \emph{System $\rightarrow$
    System}: hostname \mono{N09-Router}, múi giờ Asia/Ho\_Chi\_Minh, đồng bộ
    giờ với trình duyệt. Đồng hồ đúng là điều kiện để luật chặn theo giờ
    (Mục~\ref{subsec:nangcao}) hoạt động.
  \item \textbf{Mạng LAN và DHCP.} Nhóm giữ địa chỉ LAN mặc định
    \mono{192.168.1.1/24} và máy chủ DHCP mặc định (Hình~\ref{fig:interfaces}).
  \item \textbf{Kết nối WAN.} Cắm cáp từ mạng gia đình vào cổng WAN; giao diện
    \mono{wan} ở chế độ DHCP client tự nhận địa chỉ. Kiểm tra bằng
    \emph{Network $\rightarrow$ Diagnostics} (ping ra Internet).
  \item \textbf{Bộ phát 2,4 GHz.} Mã quốc gia VN; Mode N; kênh 6; rộng 20~MHz.
    Trong quy hoạch 20~MHz thường dùng, kênh 6 thuộc nhóm 1, 6, 11 được bố trí
    để giảm chồng lấn; lựa chọn này phù hợp với băng tần đông đúc. ESSID
    \mono{25CLC04-N09} theo quy tắc đặt tên của đề,
    Mode \emph{Access Point}, gắn vào mạng \mono{lan}.
  \item \textbf{Bộ phát 5 GHz.} Mode AC; kênh 36; rộng 80~MHz để đạt tốc độ
    cao; ESSID \mono{25CLC04-N09-5G}.
  \item \textbf{Bảo mật.} Tab \emph{Wireless Security}: WPA2-PSK, cipher
    \emph{Force CCMP (AES)}, mật khẩu dài. Bấm \emph{Save \& Apply}.
\end{enumerate}

\hinh{N09_luci_interfaces.png}{Giao diện LAN (192.168.1.1/24) và WAN (DHCP client) trên N09-Router}{fig:interfaces}
\hinh{N09_luci_wifi_24ghz.png}{Cấu hình bộ phát 2,4 GHz: mã quốc gia VN, Mode N, kênh 6, rộng 20 MHz, SSID 25CLC04-N09}{fig:wifi24}
\hinh{N09_luci_wifi_5ghz.png}{Cấu hình bộ phát 5 GHz: Mode AC, kênh 36, rộng 80 MHz, SSID 25CLC04-N09-5G}{fig:wifi5}

\subsection{Kết nối thiết bị}
\label{subsec:ketnoithietbi}
\begin{itemize}
  \item \textbf{Desktop (Linux Mint).} Rút cáp LAN để máy chỉ còn card Wi-Fi,
    chọn \mono{25CLC04-N09} trên applet mạng. Kiểm tra bằng
    \mono{nmcli device wifi} (thấy hai SSID của nhóm ở kênh 6 và 36) và
    \mono{ping -c 4 192.168.1.1}.
  \item \textbf{Laptop Windows.} Tắt \emph{Random hardware addresses}, kết nối
    vào mạng nhóm. Trang thuộc tính Wi-Fi hiển thị SSID, giao thức, kiểu bảo
    mật WPA2-Personal, băng tần, địa chỉ IPv4 và MAC; ping tới router thành
    công.
  \item \textbf{Điện thoại Android.} Kết nối mạng \mono{25CLC04-N09-5G}, chọn
    \emph{Dùng MAC của thiết bị}; trang chi tiết hiển thị địa chỉ IP, gateway
    \mono{192.168.1.1} và tần số 5~GHz.
\end{itemize}

Sau đó trạng thái toàn mạng được kiểm tra trên router
(Hình~\ref{fig:status}, Bảng~\ref{tab:leases}) và bằng ping chéo từ máy bàn
(Hình~\ref{fig:ping}, Bảng~\ref{tab:ping}).

\hinh{N09_luci_status_overview.png}{Trang Status của N09-Router: danh sách DHCPv4/DHCPv6 và trạng thái hai bộ phát}{fig:status}

\begin{table}[H]
  \centering\small
  \caption{Các thiết bị nhận địa chỉ từ N09-Router (Active DHCPv4 Leases)}
  \label{tab:leases}
  \begin{tabularx}{\linewidth}{@{}l l l L@{}}
    \toprule
    \textbf{Hostname} & \textbf{IPv4} & \textbf{MAC} & \textbf{Thiết bị}\\
    \midrule
    \mono{luckmintpc}    & 192.168.1.216 & 9C:B1:50:1F:C9:CB & Desktop Linux Mint\\
    \mono{nguyenqk-pc}   & 192.168.1.145 & C8:8A:9A:B9:1C:8E & Laptop Arch Linux (Khánh)\\
    \mono{ItzITmeehh}    & 192.168.1.199 & 24:EB:16:B6:AC:B0 & Laptop Windows (Vũ)\\
    \mono{Redmi-Note-14} & 192.168.1.115 & 42:79:95:D5:EC:D2 & Redmi Note 14 (Tuấn)\\
    (không tên)          & 192.168.1.228 & 30:56:0F:04:A9:54 & Nhiều khả năng là iPhone 11 (Khánh), kết nối không hoàn tất, không dùng trong video\\
    \bottomrule
  \end{tabularx}
\end{table}

Điện thoại Android của Vũ không có trong ảnh chụp danh sách này; minh chứng
kết nối của thiết bị này nằm ở cảnh 11 của video.

\hinh{N09_ping_crosscheck.png}{Ping chéo từ máy bàn tới các thiết bị trong mạng}{fig:ping}

\begin{table}[H]
  \centering\small
  \caption{Kết quả ping chéo từ máy bàn (192.168.1.216)}
  \label{tab:ping}
  \begin{tabularx}{\linewidth}{@{}l l c L@{}}
    \toprule
    \textbf{Đích} & \textbf{Hostname} & \textbf{Gói nhận} & \textbf{RTT min / avg / max (ms)}\\
    \midrule
    192.168.1.115 & \mono{Redmi-Note-14} & 4/4 & 7,8 / 21,1 / 54,9\\
    192.168.1.145 & \mono{nguyenqk-pc}   & 4/4 & 5,9 / 128,2 / 402,8\\
    192.168.1.199 & \mono{ItzITmeehh}    & 0/4 & \emph{Destination Host Unreachable}\\
    \bottomrule
  \end{tabularx}
\end{table}

\subsection{Các chức năng quản lý của router}
\label{subsec:nangcao}
Mỗi chức năng được thử theo cùng một quy trình: cấu hình, áp dụng, kiểm tra
trên thiết bị thử (điện thoại Android hoặc laptop), rồi trả về trạng thái ban
đầu. Bảng~\ref{tab:nangcao} ghi lại kết quả quan sát.

\begin{longtable}{@{}P{2.8cm} P{4.2cm} P{7.7cm}@{}}
  \caption{Các chức năng quản lý đã cấu hình và kết quả kiểm tra}\label{tab:nangcao}\\
  \toprule
  \textbf{Chức năng} & \textbf{Vị trí cấu hình} & \textbf{Kết quả kiểm tra}\\
  \midrule\endfirsthead
  \toprule
  \textbf{Chức năng} & \textbf{Vị trí cấu hình} & \textbf{Kết quả kiểm tra}\\
  \midrule\endhead
  \bottomrule\endfoot
  Ẩn SSID & Wireless $\rightarrow$ Edit $\rightarrow$ Hide ESSID & Mạng
  2,4~GHz biến mất khỏi danh sách quét của điện thoại thử; bỏ ẩn thì hiện
  lại\\
  Lọc địa chỉ MAC & Wireless $\rightarrow$ MAC-Filter (Allow all except listed)
  & Thiết bị có MAC trong danh sách bị ngắt và không vào lại được dù đúng mật
  khẩu; tắt lọc thì vào lại bình thường\\
  IP cố định theo MAC & DHCP and DNS $\rightarrow$ Static Leases & Máy bàn tắt
  rồi bật Wi-Fi, nhận đúng \mono{192.168.1.50} (Hình~\ref{fig:static})\\
  Chặn tên miền & DHCP and DNS $\rightarrow$ DNS Records & Trỏ
  \mono{vnexpress.net} và \mono{www.vnexpress.net} về \mono{0.0.0.0}:
  laptop Windows không mở được trang, các trang khác vẫn truy cập bình thường; xóa bản ghi thì trở lại bình thường\\
  Giới hạn giờ truy cập & Firewall $\rightarrow$ Traffic Rules (MAC nguồn,
  lan $\rightarrow$ wan, reject, Time Restrictions) & Thiết bị thử vẫn
  kết nối Wi-Fi nhưng không truy cập được web; tắt luật thì truy cập lại
  được\\
  Chuyển tiếp cổng & Firewall $\rightarrow$ Port Forwards (TCP 8080
  $\rightarrow$ 192.168.1.50:8080) & Laptop Arch Linux đặt ở mạng gia đình
  (phía WAN) truy cập được web server chạy trên máy bàn qua địa chỉ WAN của
  router; luật được xóa sau khi thử\\
  NAT & Firewall $\rightarrow$ Zones (wan: Masquerading) & Vùng \mono{wan} bật
  Masquerading, tức NAT cho toàn bộ mạng LAN\\
  Chẩn đoán & Network $\rightarrow$ Diagnostics & Traceroute tới 8.8.8.8
  đi qua router gia đình trước khi ra mạng nhà cung cấp, tức có hai lớp NAT\\
  Sao lưu cấu hình & System $\rightarrow$ Backup / Flash Firmware & Tạo được
  tệp sao lưu để phục hồi khi cần\\
\end{longtable}

\hinh[0.85\linewidth]{N09_luci_static_lease.png}{Gán IP cố định 192.168.1.50 theo MAC cho máy bàn (Static Lease)}{fig:static}

\subsection{Tổng hợp các chức năng router hỗ trợ}
Đề yêu cầu cấu hình tất cả chức năng router hỗ trợ.
Bảng~\ref{tab:tonghop} liệt kê đầy đủ các chức năng của bản OpenWrt này, kèm
trạng thái thực hiện và lý do với các chức năng không thực hiện.

\begin{table}[H]
  \centering\small
  \caption{Tổng hợp chức năng router và trạng thái thực hiện}
  \label{tab:tonghop}
  \begin{tabularx}{\linewidth}{@{}L c c L@{}}
    \toprule
    \textbf{Chức năng} & \textbf{Hỗ trợ} & \textbf{Thực hiện} & \textbf{Ghi chú}\\
    \midrule
    SSID, bảo mật WPA2 trên hai băng tần & Có & Có & Mục~\ref{subsec:cauhinh}\\
    Kênh, độ rộng kênh, chuẩn, mã quốc gia & Có & Có & Mục~\ref{subsec:cauhinh}\\
    Mật khẩu quản trị, hostname, múi giờ & Có & Có & Mục~\ref{subsec:cauhinh}\\
    LAN, DHCP, WAN (DHCP client) & Có & Có & Giữ dải LAN mặc định\\
    Ẩn SSID, lọc MAC, IP cố định & Có & Có & Bảng~\ref{tab:nangcao}\\
    Chặn tên miền, giới hạn giờ & Có & Có & Qua DNS và luật tường lửa\\
    Chuyển tiếp cổng, NAT, chẩn đoán, sao lưu & Có & Có & Bảng~\ref{tab:nangcao}\\
    Mạng khách tách riêng & Có & Không & Không đưa vào video\\
    Định tuyến tĩnh & Có & Không & Mạng chỉ có một mạng con\\
    DMZ & Không có mục riêng & -- & Tương đương chuyển tiếp mọi cổng; rủi ro cao\\
    WPS, UPnP, tên miền động (DDNS) & Không & -- & Không cài trên bản firmware
    do bộ nhớ flash chỉ 8~MB\\
    \bottomrule
  \end{tabularx}
\end{table}

% =====================================================================
\section{Thực nghiệm mô hình Ad-hoc}
\label{sec:adhoc}

\subsection{Lý do dùng máy ảo}
Các thiết bị Windows 11, Android và iOS của nhóm đều không cho người dùng tạo
mạng IBSS. Trên Windows, lệnh \mono{netsh wlan show drivers} chỉ cho biết khả
năng \emph{Hosted network}, vốn là chế độ phát Wi-Fi kiểu AP mềm (SoftAP) chứ
không phải Ad-hoc, nên không dùng được để kết luận về IBSS. Theo hướng dẫn
của đề, nhóm thực hiện trên máy ảo Ubuntu~22.04 (VMware). Máy ảo không có card
Wi-Fi thật nên nhóm dùng module nhân Linux \mono{mac80211\_hwsim} để giả lập
card Wi-Fi, mỗi card giả lập có MAC, kênh và công suất phát riêng.

\subsection{Các bước thực hiện}
\begin{enumerate}
  \item \textbf{Tạo hai card Wi-Fi ảo} và kiểm tra (Hình~\ref{fig:hwsim}):
\begin{lstlisting}
sudo modprobe mac80211_hwsim radios=2
iw dev        # xuat hien phy#0 (wlan0) va phy#1 (wlan1)
\end{lstlisting}
  \item \textbf{Tách hai card vào hai network namespace} để mỗi card là một
    ``máy'' độc lập. Nếu để chung, Linux nhận ra cả hai địa chỉ IP là của
    chính nó và chuyển gói qua loopback thay vì qua sóng ảo
    (Hình~\ref{fig:netns}):
\begin{lstlisting}
sudo ip netns add ns1
sudo ip netns add ns2
sudo iw phy phy0 set netns name ns1
sudo iw phy phy1 set netns name ns2
\end{lstlisting}
  \item \textbf{Chuyển sang chế độ IBSS và gia nhập mạng} \mono{N09-AdHoc} ở
    tần số 2437~MHz (kênh 6). Giao diện phải được tắt trước khi đổi chế độ.
    Không có DHCP nên địa chỉ IP được gán tĩnh (Hình~\ref{fig:ibsscmd}):
\begin{lstlisting}
# Node 1 (ns1); Node 2 lam tuong tu voi ns2, wlan1, 10.0.0.2/24
sudo ip netns exec ns1 ip link set wlan0 down
sudo ip netns exec ns1 iw dev wlan0 set type ibss
sudo ip netns exec ns1 ip link set wlan0 up
sudo ip netns exec ns1 iw dev wlan0 ibss join N09-AdHoc 2437
sudo ip netns exec ns1 ip addr add 10.0.0.1/24 dev wlan0
\end{lstlisting}
  \item \textbf{Kiểm tra ở tầng 802.11} bằng \mono{iw dev <card> info} và
    \mono{iw dev wlan0 station dump}; \textbf{kiểm tra ở tầng IP} bằng ping hai
    chiều giữa 10.0.0.1 và 10.0.0.2.
\end{enumerate}

\hinh[\linewidth]{N09_sodo2.pdf}{Sơ đồ mạng Ad-hoc (IBSS) hai nút trên máy ảo Ubuntu}{fig:sodo-adhoc}
\hinh[0.8\linewidth]{Picture2.png}{Hai card Wi-Fi ảo phy\#0 và phy\#1 sau khi nạp mac80211\_hwsim}{fig:hwsim}
\hinh[0.8\linewidth]{Picture3.png}{Mỗi card nằm trong một network namespace riêng}{fig:netns}
\hinh[0.85\linewidth]{Picture4.png}{Chuỗi lệnh chuyển hai nút sang chế độ IBSS (lần chạy thử với SSID demoAdhoc)}{fig:ibsscmd}

\subsection{Kết quả}
\hinh{N09_adhoc_iw_info.png}{Hai giao diện ở chế độ IBSS, SSID N09-AdHoc, kênh 6 (2437 MHz)}{fig:ibssinfo}
\hinh[0.75\linewidth]{N09_adhoc_station_dump.png}{Node 1 thấy Node 2 (02:00:00:00:01:00) ở trạng thái authorized, authenticated, associated}{fig:station}
\hinh[0.85\linewidth]{Picture6.png}{Ping từ Node 1 (10.0.0.1) tới Node 2 (10.0.0.2): 4/4 gói, mất 0\,\%}{fig:adhocping}

\begin{table}[H]
  \centering\small
  \caption{Kết quả kiểm tra mạng Ad-hoc hai nút}
  \label{tab:adhoc}
  \begin{tabularx}{\linewidth}{@{}P{3.4cm} P{3.8cm} L@{}}
    \toprule
    \textbf{Nội dung} & \textbf{Lệnh} & \textbf{Kết quả}\\
    \midrule
    Chế độ hoạt động & \mono{iw dev wlan0/wlan1 info} & \mono{type IBSS}, SSID
    \mono{N09-AdHoc}, kênh 6, rộng 20~MHz (no HT)\\
    Phát hiện nút lân cận & \mono{iw dev wlan0 station dump} & Thấy
    \mono{02:00:00:00:01:00}; authorized, authenticated, associated: yes;
    tín hiệu $-30$~dBm; chu kỳ beacon 100~TU\\
    Liên lạc tầng IP & \mono{ping -c 4} hai chiều & 4/4 gói mỗi chiều, mất 0\,\%,
    RTT khoảng 0,05--0,27~ms\\
    \bottomrule
  \end{tabularx}
\end{table}

Station dump xác nhận hai nút liên kết trực tiếp với nhau ở tầng 802.11 mà
không qua AP nào; ping thành công xác nhận liên lạc ở tầng IP. Độ trễ rất thấp
vì \mono{mac80211\_hwsim} truyền khung bằng phần mềm trong cùng nhân hệ điều
hành; kết quả chứng minh đúng nguyên lý IBSS nhưng không phản ánh suy hao,
khoảng cách hay nhiễu của Wi-Fi thật.

% =====================================================================
\section{Kết quả và thảo luận}
\label{sec:ketqua}

\subsection{Đối chiếu với yêu cầu đề}
\begin{table}[H]
  \centering\small
  \caption{Đối chiếu yêu cầu đề bài với nội dung báo cáo và video}
  \label{tab:doichieu}
  \begin{tabularx}{\linewidth}{@{}L P{6.2cm}@{}}
    \toprule
    \textbf{Yêu cầu} & \textbf{Minh chứng}\\
    \midrule
    Khái niệm WLAN & Mục~\ref{sec:khainiem}\\
    Thành phần WLAN & Mục~\ref{sec:thanhphan}\\
    Mô hình Infrastructure, Ad-hoc và giải pháp kết nối & Mục~\ref{sec:mohinh},
    \ref{sec:baomat}, \ref{sec:giaiphap}\\
    SSID theo tên nhóm (xxCLCxx-Nxx) & \mono{25CLC04-N09},
    \mono{25CLC04-N09-5G} (Bảng~\ref{tab:thongso})\\
    Cấu hình các chức năng router & Mục~\ref{subsec:cauhinh}, \ref{subsec:nangcao};
    Bảng~\ref{tab:tonghop}\\
    Kết nối Laptop, Desktop, Smartphone có minh chứng & Mục~\ref{subsec:ketnoithietbi};
    Bảng~\ref{tab:leases}, \ref{tab:ping}\\
    Ad-hoc trên máy ảo khi thiết bị không hỗ trợ & Mục~\ref{sec:adhoc}\\
    Video có thuyết minh và thấy người thực hiện & Phụ lục~\ref{app:video}\\
    Ảnh chụp minh chứng phần cấu hình & Mục~\ref{sec:infra}, \ref{sec:adhoc}\\
    Tỷ lệ đóng góp có chữ ký & Bảng B, phần đầu báo cáo\\
    \bottomrule
  \end{tabularx}
\end{table}

\subsection{Nhận xét}
\label{subsec:nhanxet}
\begin{enumerate}
  \item \textbf{Mô hình Infrastructure hoạt động đúng lý thuyết.} Router cấp
    địa chỉ DHCP cho các thiết bị trong cùng mạng con 192.168.1.0/24, hai bộ
    phát có hai BSSID khác nhau đúng như khái niệm BSS, và các thiết bị liên
    lạc được với nhau thông qua router.
  \item \textbf{Mọi trạm đều liên kết vào băng 5 GHz.} Trạng thái cuối buổi
    (Hình~\ref{fig:status}, \ref{fig:overview}) cho thấy bộ phát 5~GHz có 3
    trạm liên kết, còn SSID 2,4~GHz ở trạng thái \emph{disabled} với kiểu mã
    hóa \emph{WPA2 802.1X}. Nhiều khả năng nguyên nhân là kiểu xác thực của SSID 2,4~GHz bị
    đặt thành chế độ Enterprise trong khi mạng không có máy chủ RADIUS, khiến
    hostapd không khởi tạo được giao diện (chưa đối chiếu nhật ký hệ thống). Đây là bài học thực tế: SSID gia
    đình phải dùng chế độ Personal (PSK/SAE).
  \item \textbf{Bản ghi DHCP không chứng minh thiết bị đang trực tuyến.}
    Laptop Windows \mono{ItzITmeehh} vẫn còn trong danh sách thuê nhưng ping trả
    về \emph{Destination Host Unreachable}. Thông báo này nghĩa là máy bàn không
    nhận được phản hồi ARP, tức laptop đã rời mạng hoặc card Wi-Fi đang ngủ ở
    thời điểm đo; nếu chỉ do tường lửa Windows chặn ICMP thì kết quả sẽ là hết
    thời gian chờ (\emph{timeout}). Muốn kết luận phải đối chiếu danh sách trạm
    đang liên kết.
  \item \textbf{Địa chỉ MAC ngẫu nhiên.} MAC \mono{42:79:95:D5:EC:D2} của
    Redmi Note 14 của Tuấn có chữ số thứ hai là 2, tức đang dùng MAC ngẫu nhiên. Ví dụ
    này cho thấy vì sao phải tắt tính năng đó trước khi thử lọc MAC và chặn
    theo giờ.
  \item \textbf{Độ trễ Wi-Fi dao động mạnh.} RTT trung bình tới laptop là
    128~ms, có gói lên tới 403~ms, trong khi mạng có dây thường dưới 1~ms.
    Nhiều khả năng do cơ chế tiết kiệm năng lượng của card không dây và việc
    tranh chấp kênh theo CSMA/CA.
  \item \textbf{Mô hình Ad-hoc} chứng minh được liên kết trực tiếp giữa hai
    nút không qua AP, nhưng chỉ trên card giả lập, phù hợp với việc Ad-hoc
    hầu như không còn được hỗ trợ trên hệ điều hành phổ thông.
\end{enumerate}

\hinh[0.85\linewidth]{N09_luci_wireless_overview.png}{Wireless Overview cuối buổi: SSID 5 GHz hoạt động (WPA2-PSK), SSID 2,4 GHz bị vô hiệu do đặt nhầm WPA2-802.1X}{fig:overview}

\subsection{Hạn chế và hướng phát triển}
Video chưa trình diễn kết nối iPhone (lần kết nối thử không hoàn tất) và
mạng khách tách riêng; ở phần kết nối thiết bị, các máy chỉ được kiểm tra liên
lạc trong mạng nội bộ vì cổng WAN đã được rút ra. Nếu có thêm thời gian, nhóm sẽ: bổ sung thử nghiệm
iPhone và mạng khách; đo thông lượng bằng \mono{iperf3} trên hai băng tần;
thử chuyển vùng giữa hai AP cùng SSID (ESS); và lặp lại thử nghiệm Ad-hoc trên
hai máy Linux có card Wi-Fi thật hỗ trợ IBSS.

% =====================================================================
\section{Kết luận}
\label{sec:ketluan}
Đồ án đã hệ thống hóa khái niệm, thành phần, chuẩn kỹ thuật, bảo mật và hai mô
hình triển khai của WLAN, kèm giải pháp kết nối cho từng loại thiết bị và các
lỗi thường gặp. Về thực hành, nhóm cấu hình router TP-Link Archer C20 chạy
OpenWrt từ trạng thái mặc định thành mạng WLAN hai băng tần mang SSID
\mono{25CLC04-N09}, kết nối Desktop, Laptop và Smartphone, và thử các chức năng
quản lý như ẩn SSID, lọc MAC, IP cố định, chặn tên miền, giới hạn giờ và
chuyển tiếp cổng. Mô hình Ad-hoc được kiểm chứng trên máy ảo Linux với hai nút
IBSS liên lạc trực tiếp với nhau.

Kết quả cho thấy Infrastructure là mô hình phù hợp cho mạng gia đình nhờ có
Internet, cấp địa chỉ tự động, quản lý tập trung và bảo mật hiện đại; Ad-hoc
chỉ còn vai trò trong các mạng chuyên dụng không có hạ tầng. Các sự cố gặp
trong quá trình làm, như chọn nhầm chế độ bảo mật Enterprise hay địa chỉ MAC
ngẫu nhiên, đã được phân tích và đưa vào bảng lỗi thường gặp.

% =====================================================================
\clearpage
\phantomsection
\addcontentsline{toc}{section}{Tài liệu tham khảo}
\begin{thebibliography}{99}\small
  \bibitem{ieee80211} IEEE, \emph{IEEE Std 802.11-2020: Wireless LAN Medium
  Access Control (MAC) and Physical Layer (PHY) Specifications}, 2021.
  \bibitem{gast2005} M.~S. Gast, \emph{802.11 Wireless Networks: The Definitive
  Guide}, 2nd ed. O'Reilly Media, 2005.
  \bibitem{kurose2012} J.~F. Kurose and K.~W. Ross, \emph{Computer Networking:
  A Top-Down Approach}, 6th ed. Pearson, 2012.
  \bibitem{perahia2013} E.~Perahia and R.~Stacey, \emph{Next Generation Wireless
  LANs: 802.11n and 802.11ac}, 2nd ed. Cambridge University Press, 2013.
  \bibitem{coleman2021} D.~D. Coleman and D.~A. Westcott, \emph{CWNA Certified
  Wireless Network Administrator Study Guide}. Sybex, 2021.
  \bibitem{wpa3spec} Wi-Fi Alliance, \emph{WPA3 Specification}, 2018.
  \bibitem{krack2017} M.~Vanhoef and F.~Piessens, ``Key Reinstallation Attacks:
  Forcing Nonce Reuse in WPA2,'' in \emph{Proc. ACM CCS}, 2017.
  \bibitem{fms2001} S.~Fluhrer, I.~Mantin and A.~Shamir, ``Weaknesses in the Key
  Scheduling Algorithm of RC4,'' in \emph{Selected Areas in Cryptography},
  LNCS 2259, 2001.
  \bibitem{tt012025} Bộ Khoa học và Công nghệ, \emph{Thông tư 01/2025/TT-BKHCN
    sửa đổi, bổ sung Thông tư 08/2021/TT-BTTTT quy định danh mục thiết bị vô
  tuyến điện được miễn giấy phép}, 2025.
  \bibitem{openwrt} OpenWrt Project, \emph{OpenWrt User Guide: Wireless
  configuration, Firewall, DHCP and DNS}.
  \url{https://openwrt.org/docs/guide-user/start}
  \bibitem{hwsim} The Linux Kernel Documentation, \emph{mac80211\_hwsim}.
  \url{https://docs.kernel.org/networking/mac80211_hwsim/mac80211_hwsim.html}
  \bibitem{mswifi} Microsoft, \emph{Netsh commands for Wireless Local Area
  Network (WLAN)}. \url{https://learn.microsoft.com/windows-server/networking/technologies/netsh/netsh-wlan}
  \bibitem{android} Google, \emph{Android Help: Connect to Wi-Fi networks on
  your Android device}. \url{https://support.google.com/android}
  \bibitem{apple} Apple, \emph{Use private Wi-Fi addresses on Apple devices}.
  \url{https://support.apple.com/102509}
\end{thebibliography}

% =====================================================================
\clearpage
\appendix
\renewcommand{\thesection}{\Alph{section}}
\section{Bản đồ phân cảnh video}
\label{app:video}
Video dài 40 phút 13 giây, gồm 22 cảnh; mỗi cảnh là một clip có thuyết minh
của người trực tiếp thao tác, được dựng theo thứ tự dưới đây. Dựng video:
Nguyễn Quốc Khánh.

\begin{longtable}{@{}P{1.1cm} P{8.2cm} P{5.4cm}@{}}
  \caption{Bản đồ phân cảnh video thực hành}\label{tab:canh}\\
  \toprule
  \textbf{Cảnh} & \textbf{Nội dung} & \textbf{Thực hiện}\\
  \midrule\endfirsthead
  \toprule
  \textbf{Cảnh} & \textbf{Nội dung} & \textbf{Thực hiện}\\
  \midrule\endhead
  \bottomrule\endfoot
  \multicolumn{3}{@{}l}{\textit{Phần 0 -- Mở đầu}}\\
  01 & Giới thiệu nhóm, mục tiêu và thiết bị & Vũ, Tuấn, Khánh\\
  \multicolumn{3}{@{}l}{\textit{Phần 1 -- Giới thiệu router}}\\
  02 & Router TP-Link Archer C20 v5: mặt trước, mặt sau, nhãn, phần cứng, OpenWrt & Tuấn\\
  03 & Cách nối dây: máy quản trị cắm cổng LAN, cáp mạng nhà vào cổng WAN & Tuấn\\
  \multicolumn{3}{@{}l}{\textit{Phần 2 -- Cấu hình cơ bản}}\\
  04 & Đăng nhập LuCI lần đầu, đặt mật khẩu quản trị & Khánh\\
  05 & Đặt hostname N09-Router và múi giờ & Khánh\\
  07 & Nối WAN, kiểm tra Internet trên router & Tuấn\\
  08 & Bật Wi-Fi hai băng tần, đặt SSID, kênh, bảo mật & Tuấn\\
  \multicolumn{3}{@{}l}{\textit{Phần 3 -- Kết nối thiết bị (Infrastructure)}}\\
  09 & Desktop Linux Mint kết nối bằng Wi-Fi, \mono{nmcli}, ping & Tuấn\\
  10 & Laptop Windows: tắt MAC ngẫu nhiên, kết nối, xem thuộc tính, ping & Vũ\\
  11 & Điện thoại Android kết nối băng 5~GHz & Vũ\\
  14 & Kiểm tra toàn mạng trên router và ping chéo & Tuấn, Vũ\\
  \multicolumn{3}{@{}l}{\textit{Phần 4 -- Chức năng quản lý}}\\
  15 & Ẩn SSID & Tuấn, Vũ\\
  17 & Lọc địa chỉ MAC & Tuấn, Vũ\\
  18 & IP cố định theo MAC & Tuấn, Vũ\\
  19 & Chặn tên miền qua DNS & Tuấn\\
  20 & Giới hạn giờ truy cập & Tuấn\\
  21 & Chuyển tiếp cổng (thử từ laptop Arch ở mạng gia đình), NAT, traceroute, sao lưu cấu hình & Tuấn\\
  \multicolumn{3}{@{}l}{\textit{Phần 5 -- Mô hình Ad-hoc}}\\
  22 & Lý do dùng máy ảo & Hà Huy Hoàng\\
  23 & Tạo hai card Wi-Fi ảo, tách namespace & Hà Huy Hoàng\\
  24 & Chuyển hai nút sang IBSS, gán IP & Hà Huy Hoàng\\
  25 & Kiểm tra \mono{iw info}, \mono{station dump}, ping hai chiều & Hà Huy Hoàng\\
  \multicolumn{3}{@{}l}{\textit{Phần 6 -- Kết}}\\
  26 & Tổng kết & Vũ, Tuấn, Khánh\\
\end{longtable}

Số hiệu cảnh giữ theo kịch bản gốc; các cảnh 06, 12, 13 và 16 không được đưa
vào bản dựng cuối.

% =====================================================================
\clearpage
\section{Danh mục từ viết tắt}
\label{app:viettat}
\begin{longtable}{@{}P{2.2cm} P{6.3cm} P{6.2cm}@{}}
  \caption{Danh mục từ viết tắt}\label{tab:viettat}\\
  \toprule
  \textbf{Viết tắt} & \textbf{Tiếng Anh} & \textbf{Tiếng Việt}\\
  \midrule\endfirsthead
  \toprule
  \textbf{Viết tắt} & \textbf{Tiếng Anh} & \textbf{Tiếng Việt}\\
  \midrule\endhead
  \bottomrule\endfoot
  AES & Advanced Encryption Standard & Chuẩn mã hóa nâng cao\\
  AP & Access Point & Điểm truy cập\\
  ARP & Address Resolution Protocol & Giao thức phân giải địa chỉ\\
  BSS & Basic Service Set & Tập dịch vụ cơ bản\\
  BSSID & Basic Service Set Identifier & Định danh tập dịch vụ cơ bản\\
  CCMP & Counter Mode CBC-MAC Protocol & Giao thức mã hóa dùng AES của WPA2\\
  CSMA/CA & Carrier Sense Multiple Access with Collision Avoidance & Đa truy cập
  cảm nhận sóng mang, tránh xung đột\\
  DFS & Dynamic Frequency Selection & Chọn tần số động\\
  DHCP & Dynamic Host Configuration Protocol & Giao thức cấu hình địa chỉ động\\
  DNS & Domain Name System & Hệ thống tên miền\\
  DS & Distribution System & Hệ thống phân phối\\
  EAP & Extensible Authentication Protocol & Giao thức xác thực mở rộng\\
  ESS & Extended Service Set & Tập dịch vụ mở rộng\\
  GCMP & Galois/Counter Mode Protocol & Giao thức mã hóa Galois/Counter\\
  IBSS & Independent Basic Service Set & Tập dịch vụ cơ bản độc lập (Ad-hoc)\\
  ICMP & Internet Control Message Protocol & Giao thức bản tin điều khiển Internet\\
  IoT & Internet of Things & Internet vạn vật\\
  LLC & Logical Link Control & Điều khiển liên kết logic\\
  LuCI & Lua Configuration Interface & Giao diện quản trị web của OpenWrt\\
  MAC & Medium Access Control & Điều khiển truy cập môi trường\\
  MANET & Mobile Ad-hoc Network & Mạng di động tùy biến\\
  MIMO & Multiple Input Multiple Output & Đa đầu vào đa đầu ra\\
  MLO & Multi-Link Operation & Hoạt động đa liên kết\\
  NAT & Network Address Translation & Biên dịch địa chỉ mạng\\
  OFDM(A) & Orthogonal Frequency Division Multiplexing (Multiple Access) &
  Ghép kênh (đa truy cập) phân chia theo tần số trực giao\\
  OWE & Opportunistic Wireless Encryption & Mã hóa không dây cơ hội\\
  PHY & Physical Layer & Tầng vật lý\\
  PMF & Protected Management Frames & Bảo vệ khung quản lý\\
  PSK & Pre-Shared Key & Khóa chia sẻ trước\\
  RADIUS & Remote Authentication Dial-In User Service & Dịch vụ xác thực người
  dùng từ xa\\
  SAE & Simultaneous Authentication of Equals & Xác thực đồng thời giữa các bên
  ngang hàng\\
  SSID & Service Set Identifier & Tên mạng Wi-Fi\\
  STA & Station & Trạm, thiết bị đầu cuối\\
  TKIP & Temporal Key Integrity Protocol & Giao thức toàn vẹn khóa tạm thời\\
  TWT & Target Wake Time & Thời điểm đánh thức mục tiêu\\
  WEP & Wired Equivalent Privacy & Bảo mật tương đương mạng có dây\\
  WLAN & Wireless Local Area Network & Mạng cục bộ không dây\\
  WLC & Wireless LAN Controller & Bộ điều khiển WLAN\\
  WNIC & Wireless Network Interface Card & Card mạng không dây\\
  WPA & Wi-Fi Protected Access & Chuẩn bảo mật truy cập Wi-Fi\\
  WPAN & Wireless Personal Area Network & Mạng cá nhân không dây\\
  WPS & Wi-Fi Protected Setup & Thiết lập Wi-Fi nhanh\\
  WWAN & Wireless Wide Area Network & Mạng diện rộng không dây\\
\end{longtable}

\end{document}
